\begin{abstract}
Automated remediation is increasingly delegated to learned policies and large language models
(LLMs), yet whether a fix actually restores service, how fast, and how safely is rarely measured. We
present \sys{}, an open-source resilience lab that injects seven real faults into a three-tier web
service under open-model load, measures recovery per request, and lets remediation policies act only
through a safety shield (approved environments, an action allowlist, no repeats, an action budget,
and human approval for high-impact actions). We compare a conventional runbook, expert rules, a
LinUCB contextual bandit, gpt-4o-mini, and a tool-using gpt-4o-mini agent on training and held-out
faults, run a $2\times2$ ablation of tools and generic SRE guidance, and add an always-on agent that is
paged by SLO burn-rate alerts and heals faults it was not told about. No single action fixes every
fault: a wrong choice costs 447--1{,}128 failed requests per incident instead of 33--252. The bandit
reaches \lcBandit{} best-action accuracy (runbook \lcRunbook{}, expert rules \lcHeuristic{}) and
generalises to an unseen network partition (3/3, 0.8\,s) where the runbook fails 0/3. For the LLM,
generic guidance replaced a restart-first habit with a bypass-first habit (Redis latency
50.1\,s $\rightarrow$ 3.3\,s) but then never restarted a crashed database; tools alone added 8--12\,s
without changing decisions, and tools with guidance handled both fault types. The on-call agent
healed \ocHealed{} of \ocFaults{} unannounced faults, with medians of \ocMedDetect\,s to the alert and
\ocMedRecover\,s to healthy traffic. No unsafe action ever executed, and a Jenkins pipeline that gates
releases on resilience, alert-rule, security and on-call checks blocks a one-character bug that
passes unit and smoke tests.
\end{abstract}
